\documentclass[border=4pt]{standalone}
\usepackage{amsmath,amssymb,mathtools,xcolor,varwidth}
\definecolor{ellblue}{HTML}{1E6FE0}
\definecolor{fociorange}{HTML}{E8770A}
\definecolor{radgreen}{HTML}{00A35A}
\begin{document}
\begin{varwidth}{28em}
\large\bfseries
Kepler discovered that the planets sweep out \textcolor{ellblue}{ellipses}, not circles,      with the Sun sitting at one focus rather than at the centre. Newton now poses      the crowning question of Book I: what force law actually PRODUCES such an orbit?      The two \textcolor{fociorange}{foci $S$ and $H$} lie on the major axis, and $S$ is the force      centre while $H$ is the empty focus. Drawing the \textcolor{radgreen}{radius $SP$} from the      force to the body $P$, we measure force by the Prop. VI apparatus and will find      it varies as $1/SP^2$.
\end{varwidth}
\end{document}
